% Generated draft: requires mathematical and source review.
\documentclass[11pt]{article}
\usepackage[margin=25mm]{geometry}
\usepackage{fontspec}
\setmainfont{Latin Modern Roman}
\usepackage{amsmath,amssymb,mathtools,unicode-math}
\setmathfont{Latin Modern Math}
\usepackage{enumitem,needspace,xurl,hyperref}
\setlength{\parindent}{0pt}
\setlength{\parskip}{5pt}
\setlength{\emergencystretch}{3em}
\newcommand{\N}{\mathbb{N}}
\newcommand{\Z}{\mathbb{Z}}
\newcommand{\Q}{\mathbb{Q}}
\newcommand{\R}{\mathbb{R}}
\newcommand{\C}{\mathbb{C}}
\newcommand{\F}{\mathbb{F}}
\newcommand{\A}{\mathbb{A}}
\newcommand{\PP}{\mathbb{P}}
\newcommand{\E}{\mathbb{E}}
\newcommand{\eps}{\varepsilon}
\newcommand{\dd}{\,\mathrm{d}}
\begin{document}
\section*{UnsolvedMath ID: 20000993}
\textbf{Source problem code:} catalogue entry
\subsection{Definitions and mathematical statement}
For a finite connected graph $X$, let $J(X)=\operatorname{Div}^0(X)/\operatorname{im}L_X$ be its critical group. The source asks, through Vic Reiner, how critical-group homomorphisms associated with Galois graph coverings relate to class field theory when the deck group is abelian. It supplies no particular reciprocity statement or proposed classification.

Let $p:Y\to X$ be a connected regular unramified graph cover with finite abelian deck group $A$, $|A|=n$. Pullback sums over vertex fibers and pushforward sends each vertex to its image. The induced maps satisfy
\[
p_*p^*=n,\qquad p^*p_*=N_A:=\sum_{a\in A}a.
\]
Here pullback on critical groups is injective and pushforward surjective. Put $M=J(Y)$ and $I_AM=\langle am-m:a\in A,m\in M\rangle$. The attached aid's two integral defect identifications are
\[
J(Y)^A/p^*J(X)\simeq\widehat H^0(A,M)=M^A/N_AM,
\]
\[
\ker(p_*)/I_AM\simeq\widehat H^{-1}(A,M)=\ker N_A/I_AM.
\]
Both groups are killed by $n$. This compares ambiguous classes and a Prym coinvariant defect with Tate norm formalism, without asserting an arithmetic reciprocity theorem.

For a prime $\ell\nmid n$, write $M_\ell$ for the $\ell$-primary component. The transfer idempotent $e=n^{-1}N_A$ gives the canonical splitting
\[
J(Y)_\ell=p^*J(X)_\ell\oplus\ker(p_*)_\ell.
\]
After extension to a splitting unramified coefficient ring containing the required character values, the character idempotents $e_\chi=n^{-1}\sum_{a\in A}\chi(a)^{-1}a$ refine this to the full abelian-character decomposition. The integral defects can occur at primes dividing the covering degree, where the idempotent need not be integral.

For the cyclic cover $C_{mr}\to C_m$, $m\ge3$, $r\ge1$, the aid identifies $J(C_{mr})\simeq\Z/(mr)$, $J(C_m)\simeq\Z/m$, with pushforward reduction modulo $m$ and pullback multiplication by $r$. The deck action on $J(C_{mr})$ is trivial, so both Tate defects are $\Z/r$. The norm exact sequence
\[
0\longrightarrow\Z/r\longrightarrow\Z/(mr)\longrightarrow\Z/m\longrightarrow0
\]
splits exactly when $\gcd(m,r)=1$. Since $C_m$ admits these connected cyclic covers for every $r$, its finite critical group alone cannot classify all abelian graph covers by its quotients. The source's broader class-field comparison remains a question beyond this precise bridge and obstruction.
\subsection{Short English statement}
Relate abelian graph-cover critical groups to class-field norm ideas. The supplied aid identifies two Tate-cohomological defects and prime-to-degree splittings, while cycle covers show why a fixed finite critical group cannot serve as a complete classifying reciprocity group.
\subsection{Sources}
\begin{itemize}
\item Local UnsolvedMath registry, permanent ID $20000993$, source code catalogue entry; canonical \texttt{aim-combinatorics-notes.json}, \texttt{notes[117]}.
\item American Institute of Mathematics, \emph{Problems from the AIM Chip-Firing Workshop}, problem 28, attributed to Vic Reiner: \url{https://aimath.org/pastworkshops/chipfiringproblems.pdf}. The source cites covering-map results as [31], whose complete entry is absent; an appended page header is not part of the question. No PDF inspection is claimed here.
\item Attached research attempt 1: partial result, run and registry status partially solved, rigor proved, novelty confidence low, classification PARTIAL-PROGRESS. Its Tate identifications and cycle calculation are machine generated and not peer reviewed; independently verify before citation. No authoritative full references are supplied.
\end{itemize}
\end{document}
